% ============================================================
%  Chapter 3 — RARE-RAG: Proposed Architecture
% ============================================================
\chapter{RARE-RAG: A Risk-Aware Routed Evidence-Graph Architecture}
\label{cha:rare_rag}

% TODO: write full chapter — skeleton below
% Target: 10–12 pages

% ------------------------------------------------------------
\section{Motivation}
\label{sec:rarerag_motivation}

% TODO: 1 page
% Gap: existing architectures apply the same retrieval strategy regardless of
% query complexity or risk — inefficient and not suited for high-stakes advisory.
% Hypothesis: routing by complexity + set-wise evidence selection + grounding check
% → better faithfulness with acceptable accuracy.

% ------------------------------------------------------------
\section{Architecture Overview}
\label{sec:rarerag_overview}

% TODO: 2 pages
% Four pillars:
%   1. Risk & Complexity Router
%   2. Hybrid Evidence Retrieval (BM25 + vector, cross-encoder reranking)
%   3. Set-wise Evidence Selection (SETR-inspired)
%   4. Grounding Verifier with Abstention

% Figure placeholder:
\begin{figure}[htbp]
\centering
\begin{tikzpicture}[
  proc/.style   = {rectangle, rounded corners=5pt,
                   draw=blue!50, fill=blue!6, thick,
                   minimum width=6.2cm, minimum height=1.05cm, align=center},
  dec/.style    = {diamond, draw=orange!65, fill=orange!7, thick,
                   aspect=3.6, minimum width=5.2cm, align=center, inner sep=1pt},
  start/.style  = {ellipse, draw=blue!50, fill=blue!8, thick,
                   minimum width=2.8cm, minimum height=0.9cm, align=center},
  outbox/.style = {rectangle, rounded corners=5pt,
                   draw=green!55, fill=green!8, thick,
                   minimum width=6.2cm, minimum height=1.0cm, align=center},
  abstbox/.style= {rectangle, rounded corners=4pt,
                   draw=red!50, fill=red!6, thick,
                   minimum width=2.6cm, minimum height=0.85cm,
                   align=center, font=\small},
  arr/.style    = {-{Stealth[length=5pt]}, thick, black!65},
  lbl/.style    = {font=\footnotesize, fill=white, inner sep=1.5pt},
]

% ── main column (x=0) ──────────────────────────────────────────────
\node[start]  (query)   at (0,    0  ) {Query $q$};
\node[dec]    (router)  at (0,  -2.1 ) {\texttt{Complexity Router}};
\node[proc]   (hybrid)  at (0,  -4.6 ) {\textbf{Hybrid Retrieval}\\[2pt]
  {\footnotesize BM25\,+\,dense vector\,+\,cross-encoder reranker}};
\node[proc]   (setwise) at (0,  -7.1 ) {\textbf{Set-wise Evidence Selection}\\[2pt]
  {\footnotesize greedy complementarity-aware subset scoring}};
\node[proc]   (ground)  at (0,  -9.6 ) {\textbf{Grounding Verifier}\\[2pt]
  {\footnotesize per-claim support check against retrieved context}};
\node[proc]   (gen)     at (0, -12.1 ) {\textbf{LLM Generation}\\[2pt]
  {\footnotesize grounded response with inline citations}};
\node[outbox] (resp)    at (0, -13.9 ) {Grounded Response};

% ── abstention exit (right) ────────────────────────────────────────
\node[abstbox] (abstain) at (6.2, -9.6) {Abstention\\Response};

% ── main vertical arrows ───────────────────────────────────────────
\draw[arr] (query)   -- (router);
\draw[arr] (router)  -- node[lbl, right=3pt]{complex} (hybrid);
\draw[arr] (hybrid)  -- (setwise);
\draw[arr] (setwise) -- (ground);
\draw[arr] (ground)  -- node[lbl, right=3pt]{score $\geq\tau$} (gen);
\draw[arr] (gen)     -- (resp);

% ── abstention arrow ───────────────────────────────────────────────
\draw[arr] (ground.east) -- node[lbl, above=2pt]{score $<\tau$} (abstain.west);

% ── simple/fast bypass (left side, avoids all complex stages) ──────
% router.west → left detour → gen.west
\draw[arr] (router.west)
    -- ++(-2.0, 0)       % diamond west (-2.6) → (-4.6, -2.1)
    -- ++(0,  -10.0)     % down → (-4.6, -12.1)
    -- (gen.west)        % → gen.west (-3.1, -12.1)
    node[lbl, pos=0.07, above=2pt]{simple / fast};

\end{tikzpicture}
\caption{\rarerag{} architecture. The Complexity Router classifies each query and
  dispatches simple queries directly to LLM generation via the fast bypass (left),
  while complex or high-risk queries traverse the full four-stage pipeline.
  The Grounding Verifier emits an \emph{Abstention Response} (right) when
  per-claim evidence support falls below threshold~$\tau$.}
\label{fig:rarerag_overview}
\end{figure}

% ------------------------------------------------------------
\section{Query Complexity Router}
\label{sec:router}

% TODO: 2 pages
% Input: query text
% Classifier: lightweight LLM call with few-shot examples
% Routes: fast / standard / deep
% Design trade-off: classifier latency vs. downstream quality gain

\begin{hypothesis}[Routing benefit]
  \label{hyp:routing}
  Routing queries by complexity reduces average end-to-end latency compared to
  always applying the deepest (multi-hop + grounding) pipeline, while maintaining
  answer quality on simple queries.
\end{hypothesis}

% ------------------------------------------------------------
\section{Hybrid Evidence Retrieval}
\label{sec:hybrid_retrieval}

% TODO: 1.5 pages
% BM25 + dense vector search (alpha = 0.5 default)
% RRF fusion
% Cross-encoder reranking (BGE-reranker-v2-m3)

% ------------------------------------------------------------
\section{Set-Wise Evidence Selection}
\label{sec:setwise}

% TODO: 2 pages
% Why list-wise reranking misses inter-document complementarity
% SETR approach: score evidence subsets, not individuals
% Greedy approximate selection over top-20 candidates

% ------------------------------------------------------------
\section{Grounding Verifier with Abstention}
\label{sec:grounding_verifier}

% TODO: 2 pages
% Claim extraction from generated answer
% Per-claim grounding check against retrieved context
% Abstention: if score < threshold → "I cannot reliably answer"

\begin{hypothesis}[Faithfulness]
  \label{hyp:faithfulness}
  \rarerag{} achieves higher faithfulness than all single-path RAG architectures
  as measured by the RAGAS faithfulness metric.
\end{hypothesis}

% ------------------------------------------------------------
\section{Expected Properties}
\label{sec:rarerag_hypotheses}

% Restate all 3 hypotheses clearly
% H1: highest faithfulness — verified in Chapter 6
% H2: competitive F1 despite higher latency — verified in Chapter 6
% H3: routing reduces latency vs. always-deep — partially verified in Chapter 6
